%~Mouliné par MaN_auto v.0.23.0 2020-04-28 12:18:41
\documentclass[AHL,Unicode,longabstracts]{cedram}

\usepackage{subcaption}
%%OK
%\usepackage[rightcaption]{sidecap} 
%%OK
  
\usepackage{dsfont}
\usepackage{mathtools}


\newenvironment{step}[1]{\begin{proof}[#1]}{\let\qed\relax\end{proof}}
\newenvironment{laststep}[1]{\begin{proof}[#1]}{\end{proof}}

%\newenvironment{case}[1]{\begin{proof}[#1]}{\let\qed\relax\end{proof}}
%\newenvironment{lastcase}[1]{\begin{proof}[#1]}{\end{proof}}



\DeclareMathOperator{\union}{\bigcup}
\DeclareMathOperator{\inter}{\bigcap} 
\DeclareMathOperator{\tHt}{th}

\newcommand{\llbracket}{\mathopen{[\![}}
\newcommand{\rrbracket}{\mathclose{]\!]}}








%%MATHbb désactivés %
\newcommand{\R}{\mathbb{R}}%%remplacer \R par \Rbb
\newcommand{\Z}{\mathbb{Z}}%%remplacer \Z par \Zbb
\newcommand{\N}{\mathbb{N}}%%remplacer \N par \Nbb




%%Symboles greecs
\newcommand{\ga}{\alpha}
\newcommand{\gb}{\beta}
\newcommand{\gd}{\delta}
\newcommand{\gep}{\varepsilon}
\newcommand{\gD}{\Delta}
\newcommand{\go}{\omega}
\newcommand{\gl}{\lambda}



\newcommand{\unbf}%{\mathbf{1}}
{\mathds{1}}

\newcommand{\dd}{\mathrm{d}}
\newcommand{\tf}{\textsc{f}} %% met f en smallcap !

\renewcommand{\tilde}{\widetilde}
%% remplacer \tilde par \wtilde

\let\oldhat\hat
\renewcommand*{\hat}[1]{\mathchoice{\widehat{#1}}{\widehat{#1}}{\oldhat{#1}}{\oldhat{#1}}}


%% MATHCAL Ensuremath
\newcommand{\cA}{\ensuremath{\mathcal{A}}}
\newcommand{\cH}{\ensuremath{\mathcal{H}}}
\newcommand{\cN}{\ensuremath{\mathcal{N}}}
\newcommand{\cZ}{\ensuremath{\mathcal{Z}}}



%%MATHBF Ensuremath
\newcommand{\bP}{\ensuremath{\mathbf{P}}}
\newcommand{\bE}{\ensuremath{\mathbf{E}}}


%%MATHBB Ensuremath
\newcommand{\bbE}{\ensuremath{\mathbb{E}}}
\newcommand{\bbN}{\ensuremath{\mathbb{N}}}
\newcommand{\bbP}{\ensuremath{\mathbb{P}}}
\newcommand{\bbQ}{\ensuremath{\mathbb{Q}}}
\newcommand{\bbR}{\ensuremath{\mathbb{R}}}
\newcommand{\bbZ}{\ensuremath{\mathbb{Z}}}



%%symboles désactivés car leq et geq présents "\leqslant" et "\geqslant" = leq et geq Ams Inequalities. Inutiles.
%%\DeclareMathSymbol{\leqslant}{\mathalpha}{AMSa}{"36}
%%\DeclareMathSymbol{\geqslant}{\mathalpha}{AMSa}{"3E}
%%\DeclareMathSymbol{\eset}{\mathalpha}{AMSb}{"3F}
%%\renewcommand{\leq}{\;\leqslant\;}
%%\renewcommand{\geq}{\;\geqslant\;}




%%MATH OP STANDARDS Commandes spécialisées en indices et autres
\newcommand{\maxtwo}[2]{\max_{\substack{#1 \\ #2}}}
\newcommand{\suptwo}[2]{\sup_{\substack{#1 \\ #2}}}
\newcommand{\sumtwo}[2]{\sum_{\substack{#1 \\ #2}}}
\newcommand{\limtwo}[2]{\lim_{\substack{#1 \\ #2}}}








\newcommand\cst{\mathrm{cst}}

\newcommand*\eqsubref[1]{(\subref{#1})} % y'a probablement un réglage de subcaption qui fait ça, mais des fois c'est plus rapide de réinventer l'eau tiède que de chercher dans la documentation.

%%%%%%%%%%%%%%%%%%%%%%%%%%%%%%%%%%%%%%%%%%%%%%%%%%%%%%%%%%%%%%

\graphicspath{{./figures/}}

\newcommand*{\mk}{\mkern -1mu}
\newcommand*{\Mk}{\mkern -2mu}
\newcommand*{\mK}{\mkern 1mu}
\newcommand*{\MK}{\mkern 2mu}

\hypersetup{urlcolor=purple, linkcolor=blue, citecolor=red}

\newcommand*{\romanenumi}{\renewcommand*{\theenumi}{\roman{enumi}}}
\newcommand*{\Romanenumi}{\renewcommand*{\theenumi}{\Roman{enumi}}}
\newcommand*{\alphenumi}{\renewcommand*{\theenumi}{\alph{enumi}}}
\newcommand*{\Alphenumi}{\renewcommand*{\theenumi}{\Alph{enumi}}}
\let\oldexists\exists
\renewcommand*{\exists}{\mathrel{\oldexists}}

%%%%%%%%%%%%%%%%%%%%%%%%%%%%%%%%%%%%%%%%%%%%%%%%%%%%%%%%%%%%%%
